掷一枚硬币 100 次，正面 58 次。58\% 是这一次看到的比例，但硬币的正面概率 \(\theta\) 仍然未知。若再掷 100 次，几乎不会恰好再得到 58。要从前一次数据推知 \(\theta\)，需要回答三个逐层递进的问题：

\begin{enumerate}
\tightlist
\item
  这 100 次投掷服从怎样的生成机制？独立同分布 Bernoulli 假设意味着什么、又排除了什么？\hyperref[sec:11-Mathematical-Statistics/01-Statistical-Models/01]{``数据来自怎样的生成机制''}
\item
  在经典频率学派设定中，\(\theta\) 固定未知，而样本比例 \(\bar X\) 随样本改变。参数、统计量、估计量与估计值各处在什么位置？\hyperref[sec:11-Mathematical-Statistics/01-Statistical-Models/02]{``参数未知，而统计量随机''}
\item
  样本比例作为随机变量，在 \(\theta=0.5\) 与 \(\theta=0.6\) 下的抽样分布分别是 Binomial\((100,\theta)/100\)。这种分布如何从一次观测通向标准误、区间与检验？\hyperref[sec:11-Mathematical-Statistics/01-Statistical-Models/03]{``抽样分布连接数据与推断''}
\end{enumerate}

``总体''可以是一个有限人员名单，也可以是一个概念上的数据生成分布；``重复抽样''用于定义推断程序在假设模型下的长期行为，并不要求现实中把实验重复无限次。
